\chapter{Metodología}
\label{cap:metodologia}

\section{Tipo de trabajo}

\pendiente{Caracterizar el trabajo como desarrollo tecnológico con evaluación experimental: se construye un prototipo y se valida empíricamente su viabilidad en dispositivos móviles reales.}

\section{Etapas del trabajo}

\pendiente{Describir las etapas efectivamente seguidas (captura de datos, entrenamiento, integración en Android, pruebas), contrastándolas con el plan de trabajo de la propuesta aprobada.}

\section{Herramientas utilizadas}

\pendiente{Listar y justificar brevemente las herramientas de entrenamiento (Python, TensorFlow/Keras, MediaPipe) y de desarrollo móvil (Android, Kotlin, Jetpack Compose, CameraX, TensorFlow Lite, Text-to-Speech).}

\section{Métricas y criterios de evaluación}

\pendiente{Definir las métricas del modelo (accuracy, F1 macro y micro, matriz de confusión, consistencia estadística $\sigma$) y las métricas de rendimiento en el dispositivo (FPS efectivos, latencia de inferencia, tamaño del modelo).}

\section{Diseño de las pruebas}

\pendiente{Describir las condiciones de prueba: usuarios, iluminación, ángulo y distancia de cámara, y dispositivos utilizados.}
